\section{your first section}

In this document, i want to give you a quick overview of the possibilities of latex (a limited overview).

\subsection{citation with citavi}

When you want to input a citation by using citavi 6, you can export from there the .bib file (has to be manually refreshed after every new citation). 

Then you can citate like this: \cite{Young.2002}. There is a function in citavi, where you can copy that: \cite{Young.2002}. All the citations used get displayed at the end of the file (Quellenverzeichnis).

\subsection{handling pictures}

I use two main ways of inserting images: two pictures beside each other:

\begin{minipage}{.5\textwidth}
    \centering
    \captionsetup{type=figure} % left botm right top
    \includegraphics[clip, trim=0cm 0cm 0cm 0cm, height=5cm]{black.pdf}
    \captionof{figure}{left picture}
    \label{fig:pictureLeft}
\end{minipage}
\begin{minipage}{.5\textwidth}
    \centering
    \captionsetup{type=figure} % left botm right top
    \includegraphics[clip, trim=0cm 0cm 0cm 0cm, height=5cm]{black.pdf}
    \captionof{figure}{right picture}
    \label{fig:pictureRight}
\end{minipage}

You can trim them from all 4 sides by changing the values after trim=....

If you want to insert a reference to a picture, type this: \ref{fig:pictureLeft}.

Just one picture, you can insert like that:

\begin{figure}[H]    
    \centering          %     left botm right top
    \includegraphics[clip, trim=0cm 0cm 0cm 0cm,width=.6\textwidth]{black.pdf} %.6 stands for 0.6 x textwidth
    \caption{lonely picture}
    \label{fig:lonelyPicture}
\end{figure}

\subsection{acronyms}

When you want to use acronyms, list them in the listOfAbbreviations.tex file. As an example: in der \ac{hslu} sind Rüben. The first time you use the acronym, they get written out completely, the second time only the acronym gets used: In der \ac{hslu} hat es nicht nur Rübli sondern auch Eier.

\subsection{more references}

if you want to reference the appendix: \ref{app:firstAppendix}.

\subsection{tables}

first example for a table: (reference: \ref{tab:first table}).

\begin{table}[H]
    \centering
    \caption{first table}
    \label{tab:first table}
    \begin{tabular}{|l|l|l|l|l|l|}
        \hline
        EN & DIN & 
            \begin{tabular}{@{}c@{}}Härte \\ (HB), (HV)\end{tabular}
        &
            \begin{tabular}{@{}c@{}}Rp$_{02}$ \\ (Nmm$^{-2}$)\end{tabular}
        &
            \begin{tabular}{@{}c@{}}R$_{m}$ \\ (Nmm$^{-2})$\end{tabular}
        &  
            \begin{tabular}{@{}c@{}}E \\ (kNmm$^{-2}$)\end{tabular}
        \\ \hline
        1.4435 & X2CrNiMo18-14-3 & $\leq$ 215 & $\geq$ 200 & 500-700 & 200 \\ \hline
        1.4310 & X10CrNi18-8 & 170-250 & $\geq$ 195 & 690-900 & 195 \\ \hline
    \end{tabular}
\end{table}

Another type of table at table:

\begin{table}[H]
    \centering
    \caption{second table}
    \label{tab:second table}
    \begin{tabular}{@{}lll@{}}
    \toprule
    Bereich                 & Beschreibung                                                                  & Abbildung                     \\ \midrule
    \textbf{Netzsteuerung}  &                                                                               &                               \\
                            &                                                                               &                               \\
    Form                    & Viereckform                                                                   & \ref{fig:lonelyPicture}       \\
                            &                                                                               &                               \\
    Ansatzfunktion          & Lineare Ansatzfunktion                                                        & -                             \\
                            &                                                                               &                               \\
    Verfeinerung            & Lokale Verfeinerung (Löcher Zentrum)                                          & \ref{fig:pictureLeft}         \\
                            &                                                                               &                               \\
    Elementgrösse           & Global max.: 0.02 mm                                                          & -                            \\
                            &                                                                               &                               \\
    Konvergenzanalyse       & Maximale Abweichung 1\%                                                       & -                             \\ \midrule
    \textbf{Statistik}      &                                                                               &                               \\
                            &                                                                               &                               \\
    Anzahl Elemente         & 6908                                                                          & \ref{fig:pictureRight}        \\
                            &                                                                               &                               \\
    Anzahl Knoten           & 7449                                                                          & -                             \\ \bottomrule
    \end{tabular}
\end{table}

\newpage

\subsection{formulas}

geg: d$_{innen}$ = 48.53 mm, F$_{max}$ = 1500 N\newline
ges: p$_{max}$

\begin{equation}
    p_{max}=\dfrac{F_{max}}{A}=\dfrac{1500N}{\dfrac{\pi \cdot (49.53mm)^{2}}{4}}=0.7785Nmm^{2}=7.79bar
    \label{eq:firstEquation}
\end{equation}

reference: \ref{eq:firstEquation}

You can either write an equation between two \$ signs: $in text equation \pi$ or as shown above in equation \ref{eq:firstEquation}.

For a new paragraph, let a line empty between text.

\subsection{enumerations and itemize}

Two different ways: first:

\begin{enumerate}
    \item two
    \item three
    \item five
\end{enumerate}

second way:

\begin{itemize}
    \item you
    \item Arten
    \item idk
\end{itemize}

\subsection{general information}

All the pictures you want to use put in the -images- folder. New .tex files put in the -src- folder and PDFs and stuff for the appendix put in the folder anhang-.